Composite:  (f \circ g)(x) = f(g(x))
Evaluate inner function first, then substitute into outer.

Example:  f(x) = 20x+4,\; g(x) = 24x-10
g\!\left(-\frac{1}{4}\right) = 24\!\left(-\frac{1}{4}\right) - 10 = -6 - 10 = -16
f(-16) = 20(-16) + 4 = -316

Inverse: swap  x  and  y, solve for  y.

f(x) = \frac{3x+2}{2x-3} \rightarrow y = \frac{3x+2}{2x-3}
\rightarrow x(2y-3) = 3y+2
\rightarrow 2xy - 3x = 3y + 2
\rightarrow y(2x-3) = 3x+2
\rightarrow f^{-1}(x) = \frac{3x+2}{2x-3}

Note: this function is its own inverse (involution).

f^{-1}(v) = w \rightarrow  f(w) = v \quad (use this shortcut on competition)

f\!\left(f^{-1}(x)\right) = x \qquad f^{-1}\!\left(f(x)\right) = x

Operations on functions:
(f + g)(x) = f(x) + g(x)
(fg)(x) = f(x) \cdot g(x)
\left(\frac{f}{g}\right)(x) = \frac{f(x)}{g(x)},\quad g(x) \ne 0

Domain: all  x  for which  f(x)  is defined.
Range: all output values of  f(x).

For  f = \{(3,5),(5,7),(2,3)\}: domain = \{2,3,5\}, sum = 10
